\section{Theoretical analysis}
\label{sec:theory}

Each source predictor can admit different numerical representations.
We first characterize
the information about sharing that survives positive neuron rescaling.
We then give conditions for recovering the groups, describe the
remaining ambiguities and relate candidate selection to transfer.
All proofs and additional results are in
Appendix~\ref{app:theory-proofs}.

\subsection{Sharing information after removing neuron scales}
\label{sec:theory-observations}

Consider the ReLU layer of Section~\ref{sec:method-relu-instance}.
Throughout the structural analysis, fix a common set of active
positions $\mathcal E\subseteq\{1,\ldots,D\}\times\{1,\ldots,H\}$
and a correct correspondence of hidden neurons across tasks.
Write $E:=|\mathcal E|$, enumerate these positions by $p=1,\ldots,E$,
and let $h(p)$ denote the neuron containing position $p$.
Every one of the $H$ neurons is assumed to contain an active position.

To include the paired codebook used in the experiments, let $m$ count
the used independent code values, with every group nonempty.
Describe an indexator $\mathbf I$ on $\mathcal E$ by a group assignment
$j_{\mathbf I}(p)\in\{1,\ldots,m\}$ and a fixed sign
$\sigma_{\mathbf I,p}\in\{-1,1\}$.
Ordinary equality sharing is the case
$\sigma_{\mathbf I,p}=1$.
Enumerate the $T$ source tasks by $\tau=1,\ldots,T$.
For each task, let
$\mathbf c_\tau\in(\mathbb R\setminus\{0\})^m$ be its code values
and let $\mathbf s_\tau\in(0,\infty)^H$ be its neuron scales.
Let $\mathbf r_\tau=(r_{\tau,p})_{p=1}^E$ collect the active entries
of the aligned source directions.
In the exact model, these entries satisfy
\begin{equation}
    r_{\tau,p}
    =\frac{\sigma_{\mathbf I,p}c_{\tau,j_{\mathbf I}(p)}}
           {s_{\tau,h(p)}},
    \qquad p\in\{1,\ldots,E\}.
    \label{eq:theory-observation}
\end{equation}
Positive scales express the ReLU symmetry described in
Section~\ref{sec:method-relu-instance}
\cite{stock2022embedding}. The assumptions here concern an exact
representation of the source solutions; errors in these directions
are considered in Appendix~\ref{app:theory-noise}.

Define the group-membership matrix
$\mathbf U_{\mathbf I}\in\{0,1\}^{E\times m}$ and the
neuron-membership matrix $\mathbf N\in\{0,1\}^{E\times H}$ by
\begin{equation}
    (\mathbf U_{\mathbf I})_{pj}
        :=\mathbf1\{j_{\mathbf I}(p)=j\},
    \qquad
    N_{ph}:=\mathbf1\{h(p)=h\}.
    \label{eq:theory-membership}
\end{equation}
Let $\boldsymbol\ell_\tau:=\log|\mathbf r_\tau|$,
$\mathbf q_\tau:=\log|\mathbf c_\tau|$ and
$\boldsymbol\alpha_\tau:=\log\mathbf s_\tau$,
with logarithms and absolute values applied coordinatewise.
Equation~\eqref{eq:theory-observation} gives
$\boldsymbol\ell_\tau=\mathbf U_{\mathbf I}\mathbf q_\tau
-\mathbf N\boldsymbol\alpha_\tau$.
The matrix
\begin{equation}
    \boldsymbol\Pi
    :=\mathbf{Id}_E-\mathbf N
        (\mathbf N^{\mathsf T}\mathbf N)^{-1}\mathbf N^{\mathsf T}
    \label{eq:theory-centering}
\end{equation}
subtracts the mean within each neuron.
Define the centered membership matrix
$\mathbf M_{\mathbf I}:=\boldsymbol\Pi\mathbf U_{\mathbf I}$
and the centered observations
$\mathbf y_\tau:=\boldsymbol\Pi\boldsymbol\ell_\tau$.
Since $\boldsymbol\Pi\mathbf N=\mathbf0$, they satisfy
\begin{equation}
    \mathbf y_\tau=\mathbf M_{\mathbf I}\mathbf q_\tau.
    \label{eq:theory-contrasts}
\end{equation}

\begin{definition}[Contrast space]
    The contrast space of $\mathbf I$ is
    $\mathcal S_{\mathbf I}:=\operatorname{range}(\mathbf M_{\mathbf I})$.
    It contains all centered log-magnitude observations obtainable by
    varying the nonzero code values and positive neuron scales.
\end{definition}

Thus the contrast space records the observable variation of weight
magnitudes after neuron scales have been removed.
To describe its dimension, define the \emph{calibration graph}
$G_{\mathbf I}$: its vertices are the $m$ codes and $H$ neurons,
and each active position $p$ is an edge joining
$j_{\mathbf I}(p)$ to $h(p)$. Repeated occurrences of a code in
one neuron give parallel edges. Let $g_{\mathbf I}$ be the number
of connected components of this graph.

\begin{proposition}[Observable dimension and scale freedom]
    \label{prop:theory-dimension}
    Under \eqref{eq:theory-observation},
    \begin{equation}
        \dim\mathcal S_{\mathbf I}=m-g_{\mathbf I},
        \qquad
        \operatorname{rank}[\mathbf U_{\mathbf I},-\mathbf N]
            =m+H-g_{\mathbf I}.
        \label{eq:theory-dimension}
    \end{equation}
    Each component permits one common positive rescaling of its
    code values and neuron scales. Fixing one neuron scale in each
    component removes this freedom in an exact calibration.
\end{proposition}

This explains the component-wise normalization in
Appendix~\ref{app:relu-normalization}.
The graph has $E-m-H+g_{\mathbf I}$ independent cycles.
They supply the consistency relations between observed magnitudes:
products along alternating sides of a cycle agree. These are the
classical rank-one cycle relations, which become linear after taking
logarithms \cite[Theorem~34 and Remark~35]{kiraly2015algebraic}.
On a forest, any positive edge magnitudes admit such a factorization.
Repeated tasks then provide additional magnitude constraints when the
shared grouping creates cycles.

\subsection{When the sharing groups can be recovered}
\label{sec:theory-recovery}

The next question is whether the contrast space determines the group
assignment itself. Suppose this space is known, and write
$r:=\dim\mathcal S_{\mathbf I}$.
Choose an orthonormal basis
$\mathbf Q\in\mathbb R^{E\times r}$ for it.
For two positions $p,q$ in the same neuron, let
$\mathbf e_p,\mathbf e_q\in\mathbb R^E$ be their standard basis
vectors and define the ordered-pair signature
\begin{equation}
    \boldsymbol\delta_{pq}
        :=(\mathbf e_p-\mathbf e_q)^{\mathsf T}\mathbf Q,
    \qquad h(p)=h(q).
    \label{eq:theory-signature}
\end{equation}
A zero signature identifies equal codes within a neuron.
Equal nonzero signatures for pairs $(p,q)$ and $(u,v)$ identify the
same ordered pair of codes: the first positions have one common code,
and the second positions have another.
The resulting recovery rule merges these positions and takes the
transitive closure of the merges.

For each code $j$, define a graph $W_{\mathbf I}(j)$ whose vertices
are the neurons containing $j$. Two vertices are adjacent when the
neurons share another code distinct from $j$.
A one-vertex graph is considered connected.
These graphs describe whether repeated pairs link all occurrences
of a code.

\begin{theorem}[Recovery from repeated code pairs]
    \label{thm:theory-pair-recovery}
    If every $W_{\mathbf I}(j)$ is connected, the signature rule
    recovers all groups of $\mathbf I$.
    Consequently, any grouping $\mathbf J$ on the same active
    positions and fixed neurons, with $m$ used codes and
    $\mathcal S_{\mathbf J}=\mathcal S_{\mathbf I}$, has the same
    groups up to renaming.
\end{theorem}

For example, the condition holds if one neuron contains every code
and each other neuron contains at least two distinct codes.
If some $W_{\mathbf I}(j)$ is disconnected, the same rule recovers
correct partial groups, with one part for each component of that graph.
Longer cycles can determine how such parts join;
Appendix~\ref{app:theory-long-cycles} gives an alternative
uniqueness condition using all cycles.

To relate this result to multiple source tasks, let
$\mathbf Z:=[\mathbf y_1,\ldots,\mathbf y_T]\in\mathbb R^{E\times T}$
collect their centered observations.

\begin{corollary}[Identification from a sufficiently diverse source set]
    \label{cor:theory-task-diversity}
    Let $\mathcal C$ be a finite set of candidate groupings on the
    fixed active positions and neurons, containing $\mathbf I$.
    Suppose each candidate uses $m$ codes and has contrast dimension
    $r$, every $W_{\mathbf I}(j)$ is connected, and
    $\operatorname{rank}\mathbf Z=r$.
    Every candidate in $\mathcal C$ that explains all observed
    magnitudes has the same groups as $\mathbf I$, up to renaming.
\end{corollary}

The rank condition means that the source tasks reveal every direction
of variation in the true contrast space. It is sufficient for learning
the entire space; a finite set of different candidate spaces can also
be separated by fewer observations.
Under perturbations, recovery additionally depends on the error in
the estimated space. Appendix~\ref{app:theory-noise} gives an explicit
error threshold and its dependence on small weights and task diversity.

\subsection{Remaining ambiguities and the information in signs}
\label{sec:theory-ambiguity}

The repeated-pair condition is sufficient for uniqueness.
For a general calibration graph, the remaining candidates can be
characterized by its cycle constraints.
Orient every edge of $G_{\mathbf I}$ from its code to its neuron,
choose a spanning forest, and consider the cycle formed by each
edge outside that forest.
Let $b_{\mathbf I}:=E-m-H+g_{\mathbf I}$.
Define the cycle matrix
$\mathbf R_{\mathbf I}\in\{-1,0,1\}^{b_{\mathbf I}\times E}$
by giving each cycle one row: its edge coefficient is $+1$ or $-1$
according to agreement with the traversal direction, and zero
outside the cycle. These coefficients specify edge orientation;
the weight signs remain the separate quantities
$\sigma_{\mathbf I,p}$.

For another grouping $\mathbf J$, let $m_{\mathbf J}$ be its number
of used codes and $g_{\mathbf J}$ the number of components of
$G_{\mathbf J}$. Its membership matrix
$\mathbf U_{\mathbf J}$ is defined as in
\eqref{eq:theory-membership}.

\begin{theorem}[Exact criterion for magnitude equivalence]
    \label{thm:theory-cycle-equivalence}
    On the fixed active positions and neurons,
    \begin{equation}
        \mathcal S_{\mathbf J}=\mathcal S_{\mathbf I}
        \quad\Longleftrightarrow\quad
        \left\{
        \begin{aligned}
            \mathbf R_{\mathbf I}\mathbf U_{\mathbf J}&=\mathbf0,\\
            m_{\mathbf J}-g_{\mathbf J}&=m-g_{\mathbf I}.
        \end{aligned}
        \right.
        \label{eq:theory-cycle-equivalence}
    \end{equation}
\end{theorem}

The first condition requires each new group to occur equally often
on the positive and negative sides of every fundamental cycle.
The second preserves the dimension of observable variation.
Together they describe all partitions with the same contrast space.
Their enumeration can be exponential in the number of positions.
Appendix~\ref{app:theory-counterexample} gives distinct groupings with
the same number of codes and the same contrast space.

Signs provide complementary information because positive neuron
rescaling preserves them. For each code $j$, define its sign profile
$\boldsymbol\epsilon_j:=(\operatorname{sign}c_{1j},\ldots,
\operatorname{sign}c_{Tj})\in\{-1,1\}^T$.
The observed profile of position $p$ is
$\mathbf b_p:=(\operatorname{sign}r_{1p},\ldots,
\operatorname{sign}r_{Tp})$.
Equation~\eqref{eq:theory-observation} gives
$\mathbf b_p=\sigma_{\mathbf I,p}
\boldsymbol\epsilon_{j_{\mathbf I}(p)}$.

\begin{proposition}[Recovery by sign profiles]
    \label{prop:theory-signs}
    If the code profiles are pairwise distinct up to a common sign,
    grouping the observed profiles up to sign recovers the groups.
    Comparisons with one representative per group also recover the
    relative signs of its positions.
    Among representations with $m$ used codes, the signed indexator
    is determined up to code renaming and a sign change of each code
    together with its assignments.
\end{proposition}

This result applies independently of graph connectivity and magnitude
calibration. Its usefulness depends on sign variation across tasks;
codes that keep the same sign can have identical profiles.
Finite-task collision bounds and a noisy version are given in
Appendix~\ref{app:theory-sign-details}.

These results characterize aligned weight representations.
Connecting them to the predictor classes of
Section~\ref{sec:problem-formulation} requires the corresponding
network identifiability assumptions \cite{stock2022embedding}.
The example in Appendix~\ref{app:theory-counterexample} shows
explicitly how magnitude-equivalent groupings can define different
signed ReLU function classes. Accordingly, a candidate bank can
retain several groupings within one contrast space and use
functional evaluation to distinguish them.

\subsection{Candidate selection and transfer to new tasks}
\label{sec:theory-transfer}

Structural identifiability concerns which constraints the source
representations reveal. Their usefulness on new tasks also depends
on numerical fitting. To express this dependence, fix a support size
$n$, a training budget $b$ and a numerical adaptation rule
$\mathcal A_{\mathbf I,b}$ for each admissible indexator.
Given a support sample $S$ of $n$ examples and independent
initialization randomness $\zeta$, the rule returns a predictor
$\mathcal A_{\mathbf I,b}(S,\zeta)$ after fitting all numbers anew.
Define its adaptation risk by
\begin{equation}
    \mathcal R_{n,b}^{\mathrm{adapt}}(\mathbf I)
    :=\mathbb E_{\tau,S,\zeta}
        \mathbb E_{(\mathbf x,y)\sim P_\tau}
        \ell\bigl(\mathcal A_{\mathbf I,b}(S,\zeta)(\mathbf x),y\bigr),
    \label{eq:theory-adaptation-risk}
\end{equation}
where $\tau\sim p(\tau)$ and $S$ is sampled conditionally on $\tau$
as in Section~\ref{sec:tasks}.

Consider selection on independent tasks after fixing a nonempty
finite bank $\mathcal B$ and all adaptation rules.
Let $\mathcal H$ contain $T_{\mathcal H}$ independent tasks drawn
from $p(\tau)$, independently of bank construction.
For each $\tau\in\mathcal H$, draw its support sample of size $n$,
an independent query sample
and an independent initialization $\zeta_\tau$.
The empirical selection risk is
\begin{equation}
    \widehat{\mathcal R}_{\mathcal H}(\mathbf I)
    :=\frac1{T_{\mathcal H}}\sum_{\tau\in\mathcal H}
        \mathcal L_{S_\tau^{\mathrm q}}
        \bigl(\mathcal A_{\mathbf I,b}
            (S_\tau^{\mathrm s},\zeta_\tau)\bigr).
    \label{eq:theory-meta-selection}
\end{equation}
Let $\widehat{\mathbf I}_{\mathcal H}$ minimize this quantity over
$\mathcal B$, and let $\delta\in(0,1)$ be a failure probability.

\begin{theorem}[Selection on independent tasks]
    \label{thm:theory-transfer}
    Suppose the loss lies in $[0,L]$ for a fixed $L>0$.
    With probability at least $1-\delta$,
    \begin{equation}
        \mathcal R_{n,b}^{\mathrm{adapt}}
            (\widehat{\mathbf I}_{\mathcal H})
        \le
        \min_{\mathbf I\in\mathcal B}
            \mathcal R_{n,b}^{\mathrm{adapt}}(\mathbf I)
        +2L\sqrt{\frac{\log(2|\mathcal B|/\delta)}
                         {2T_{\mathcal H}}}.
        \label{eq:theory-transfer-bound}
    \end{equation}
    This is a finite-family consequence of Hoeffding's inequality
    \cite{hoeffding1963probability} applied at the task level.
\end{theorem}

The first term measures the quality of the best candidate after
the prescribed numerical adaptation. The second measures the cost
of selecting it from independent tasks.
Coverage of the bank is a separate requirement: a missing useful
structure contributes an approximation error even with arbitrarily
accurate selection.
For squared error and cross-entropy, applying this bounded-loss
statement requires a prespecified bounded loss or an appropriate
tail analysis.

The source structure-selection samples in
Section~\ref{sec:method-selection} provide new examples of tasks
already used to construct the bank.
Their conditional guarantee concerns those source tasks
(Appendix~\ref{app:theory-selection}).
Theorem~\ref{thm:theory-transfer} describes the additional
task-independence condition for a transfer guarantee.
Appendix~\ref{app:theory-adaptation} further separates approximation
error, estimation variance and numerical conditioning in simple
models. Together, these distinctions explain why exact recovery,
prediction error and training time are evaluated separately.
